\documentclass[11pt,a4paper]{article}

% ============================================
% PACOTES
% ============================================
\usepackage[utf8]{inputenc}
\usepackage[T1]{fontenc}
\usepackage[english]{babel}
\usepackage{amsmath,amssymb,amsthm}
\usepackage{mathtools}
\usepackage{geometry}
\usepackage{graphicx}
\usepackage{hyperref}
\usepackage{cleveref}
\usepackage{booktabs}
\usepackage{enumitem}

\geometry{margin=1in}

% ============================================
% TEOREMAS E DEFINIÇÕES
% ============================================
\theoremstyle{definition}
\newtheorem{axiom}{Axiom}
\newtheorem{definition}{Definition}[section]
\newtheorem{theorem}{Theorem}[section]
\newtheorem{corollary}[theorem]{Corollary}
\newtheorem{proposition}[theorem]{Proposition}

\theoremstyle{remark}
\newtheorem{remark}{Remark}[section]
\newtheorem{example}{Example}[section]

% ============================================
% COMANDOS PERSONALIZADOS
% ============================================
\newcommand{\NS}{\mathcal{N}_s}
\newcommand{\CA}{\mathcal{A}}
\newcommand{\CO}{\Omega}
\newcommand{\Dkl}{D_{\mathrm{KL}}}
\newcommand{\Djsd}{D_{\mathrm{JSD}}}
\newcommand{\Graph}{\mathcal{G}}
\newcommand{\Yharim}{\Upsilon}
\newcommand{\Lap}{\mathcal{L}}
\newcommand{\vsilence}{v_{\mathrm{s}}}

\hypersetup{
    colorlinks=true,
    linkcolor=blue,
    citecolor=blue,
    urlcolor=blue
}

% ============================================
% METADADOS
% ============================================
\title{
    \textbf{Informational Tension Theory} \\
    \Large A Framework for Entropic Detection of Latent Structures \\[0.5em]
    \normalsize Core Paper v3.0 (With Temporal Dynamics)
}

\author{
    Matheus Grego\footnote{Independent Researcher. Mathematical formalization assisted by Claude (Anthropic) under author supervision. Full paper available with complete proofs and extended validation.}
}

\date{January 2026}

\begin{document}

\maketitle

% ============================================
% ABSTRACT
% ============================================
\begin{abstract}
We present a mathematical framework for detecting \emph{latent entities} in complex systems through entropic deformation analysis. The central thesis: suppression of information in densely connected systems generates measurable structural anomalies in adjacent nodes. We formalize \emph{Informational Tension} ($\tau$) as a divergence-based local measure and \emph{Concealment Cost} ($\Omega$) as a global measure scaling superlinearly with connectivity. 

\textbf{Empirical Validation.} Cross-domain validation on software ecosystems (GitHub) and ecological networks confirms the framework's predictions. The central finding is the \emph{Sniper Effect}: a low-degree node (developer ``Marak,'' $d=2$) exhibited tension $\tau = 0.31$ exceeding detection thresholds, despite being topologically invisible to degree-based methods. We introduce the \emph{Yharim Limit} ($\Yharim$), a fundamental signal-to-noise threshold defining the boundary of detectability.

\vspace{0.5em}
\noindent\textbf{Keywords:} Information Theory, Graph Anomaly Detection, Latent Variable Inference, Network Science
\end{abstract}

% ============================================
% 1. INTRODUCTION
% ============================================
\section{Introduction}

Consider a network representing relationships between entities---social connections, biological interactions, financial transactions, or software dependencies. When an entity is \emph{suppressed} (removed, hidden, or censored), the resulting ``hole'' in the network creates structural anomalies detectable through careful analysis.

This paper formalizes the intuition that \emph{absence leaves traces}. We develop Informational Tension Theory (ITT), a mathematical framework where:
\begin{enumerate}[noitemsep]
    \item Suppression creates measurable \emph{tension} in adjacent nodes
    \item The \emph{cost} of concealment scales superlinearly with connectivity
    \item A fundamental \emph{detectability limit} (the Yharim Limit) separates signal from noise
\end{enumerate}

The framework is \emph{domain-agnostic}: the same mathematics governs species extinction, account banning, and financial fraud detection.

\subsection{Motivating Example: The faker.js Incident}

In January 2022, developer ``Marak Squires'' deliberately corrupted his widely-used npm packages \texttt{faker.js} and \texttt{colors.js}, affecting thousands of dependent projects. He was subsequently banned from GitHub. 

Traditional network analysis would classify Marak as a low-degree node ($d=2$ in the sampled co-contribution network)---theoretically undetectable. Yet our framework reveals that Marak's \emph{informational tension} ($\tau = 0.31$) exceeded detection thresholds, exposing the suppression through the anomalous behavior of his former collaborators.

This is the \emph{Sniper Effect}: a low-degree node whose removal creates disproportionate structural deformation, detectable through entropic analysis rather than topological metrics.

\subsection{Contributions}

This paper makes the following contributions:

\begin{enumerate}
    \item \textbf{Axiomatic Foundation}: Three axioms (Conservation, Deformation, Structural Cost) grounding the theory in information-theoretic principles
    
    \item \textbf{Formal Definitions}: Informational Tension ($\tau$) and Concealment Cost ($\Omega$) with rigorous mathematical characterization
    
    \item \textbf{The Yharim Limit}: A fundamental detectability threshold $\Yharim = \sqrt{2\log(1/\alpha)}$ defining the boundary between signal and noise
    
    \item \textbf{The Sniper Effect}: Theoretical prediction and empirical validation that low-degree nodes can exhibit high tension
    
    \item \textbf{Cross-Domain Validation}: Empirical confirmation across software ecosystems and ecological networks
\end{enumerate}

% ============================================
% 2. AXIOMATIC FOUNDATIONS
% ============================================
\section{Axiomatic Foundations}

\subsection{Primitive Concepts}

\begin{definition}[Information Graph]
An \emph{information graph} is a tuple $\Graph = (V, E, \phi)$ where $V$ is a set of entities, $E \subseteq V \times V$ represents relationships, and $\phi: V \cup E \to \mathcal{F}$ assigns observable features.
\end{definition}

\begin{definition}[Silence Node]
A \emph{silence node} $\NS \in V^* \setminus V$ is an entity present in the true graph $\Graph^*$ but absent from the observed graph $\Graph$. Its \emph{adjacency layer} is $\CA(\NS) = \{v \in V : (v, \NS) \in E^*\}$.
\end{definition}

\begin{definition}[Observed vs True Graph]
Given a true graph $\Graph^* = (V^*, E^*, \phi^*)$, the \emph{observed graph} $\Graph = (V, E, \phi)$ is the accessible subgraph where $V \subseteq V^*$. The discrepancy between $\Graph^*$ and $\Graph$ encodes the information about suppressed entities.
\end{definition}

\subsection{The Three Axioms}

We now state the foundational axioms of Informational Tension Theory.

\begin{axiom}[Conservation of Relational Structure]
\label{ax:conservation}
The local structure around each node $v$ follows a probability distribution $P_{\Graph}(\cdot | v)$ determined by the generative model. Deviations require external intervention.
\end{axiom}

This axiom establishes that information graphs follow statistical regularities. Whether representing social ties, biological interactions, or financial flows, local neighborhoods exhibit predictable structural properties (degree distribution, clustering, motifs).

\begin{axiom}[Deformation Principle]
\label{ax:deformation}
Removal of node $\NS$ induces measurable \emph{deformation} in structural properties of nodes in $\CA(\NS)$, proportional to prior connection strength.
\end{axiom}

The deformation manifests as: missing edges (the direct connections to $\NS$), broken triadic closure (paths through $\NS$ are severed), and information flow disruption (signals that would propagate through $\NS$ are blocked).

\begin{axiom}[Structural Cost of Concealment]
\label{ax:structural-cost}
Maintaining suppression requires continuous structural intervention, manifesting as increased entropy in $\CA(\NS)$ detectable as statistical anomalies.
\end{axiom}

\begin{remark}[Mathematical Isomorphism]
Axiom 3 has structure \emph{isomorphic} to thermodynamic systems---maintaining an information differential requires continuous effort, analogous to maintaining a temperature differential. This is a mathematical correspondence (the governing equations have identical form), not an ontological claim that networks are thermodynamic systems.
\end{remark}

% ============================================
% 3. CORE DEFINITIONS
% ============================================
\section{Core Definitions}

\subsection{Informational Tension}

The central quantity of our framework measures local structural anomaly.

\begin{definition}[Local Informational Tension]
\label{def:tension}
For node $v_i$ in observed graph $\Graph$, the \emph{informational tension} is:
\begin{equation}
    \boxed{
    \tau(v_i) = \Dkl\Big( P_{\text{obs}}(\mathcal{N}(v_i)) \;\Big\|\; P_{\text{exp}}(\mathcal{N}(v_i)) \Big)
    }
\end{equation}
where $P_{\text{obs}}$ is the observed neighborhood distribution and $P_{\text{exp}}$ is the expected distribution under the generative model.
\end{definition}

The KL divergence measures the ``surprise'' of observing the actual neighborhood when expecting the model-predicted structure. High $\tau(v_i)$ indicates anomalous local structure---the node behaves as if something is missing.

\begin{theorem}[Tension-Suppression Correspondence]
\label{thm:correspondence}
Let $\NS$ be a silence node and $v_i \in \CA(\NS)$. Under mild regularity conditions:
\begin{equation}
    \mathbb{E}[\tau(v_i) | \NS \text{ suppressed}] > \mathbb{E}[\tau(v_i) | \NS \text{ never existed}]
\end{equation}
with the difference monotonically increasing with edge weight $w(\NS, v_i)$.
\end{theorem}

\begin{proof}[Proof Sketch]
If $\NS$ never existed, $P_{\text{obs}} \approx P_{\text{exp}}$ and $\tau \approx 0$. If $\NS$ was suppressed, the removal creates missing edges, broken paths, and clustering perturbations, causing $P_{\text{obs}}$ to deviate from $P_{\text{exp}}$.
\end{proof}

\begin{remark}[Robust Implementation]
\label{rem:robust}
In practice, we use \emph{Jensen-Shannon Divergence} (JSD) instead of KL divergence:
\begin{equation}
    \text{JSD}(P, Q) = \frac{1}{2}\Dkl(P \| M) + \frac{1}{2}\Dkl(Q \| M), \quad M = \frac{P+Q}{2}
\end{equation}
with $\text{JSD} \in [0, 1]$. This handles disjoint supports gracefully. For thresholding, we use \emph{Median Absolute Deviation} (MAD) instead of standard deviation for robustness to power-law outliers common in real networks.
\end{remark}

\subsection{Concealment Cost}

Having established that suppression creates local tension, we define a global measure.

\begin{definition}[Concealment Cost]
\label{def:concealment}
The \emph{concealment cost} of silence node $\NS$ is:
\begin{equation}
    \boxed{
    \CO(\NS) = \sum_{k=1}^{K} \sum_{v_i \in \CA^{(k)}(\NS)} \tau(v_i) \cdot \omega(k)
    }
\end{equation}
where $\omega(k) = e^{-\lambda k}$ is an exponential decay function, $\CA^{(k)}$ is the $k$-th order adjacency layer, and $\lambda > 0$ is the tension attenuation constant.
\end{definition}

\begin{theorem}[Superlinear Scaling]
\label{thm:scaling}
For silence node $\NS$ with degree $d_{\NS}$:
\begin{equation}
    \CO(\NS) = \Theta\left( d_{\NS}^{1+\alpha} \right)
\end{equation}
where $\alpha > 0$ depends on clustering coefficient. In highly clustered graphs, $\alpha \to 1$; in tree-like graphs, $\alpha \to 0$.
\end{theorem}

\begin{proof}[Proof Sketch]
Removing a high-degree node breaks not only direct connections but also paths \emph{between} neighbors. The number of broken neighbor-neighbor paths is $\binom{d_{\NS}}{2} \cdot C \sim d_{\NS}^2$. Combined with the linear term from direct neighbors, we get:
\begin{equation}
    \CO(\NS) \sim d_{\NS} \cdot \tau_{\text{direct}} + \binom{d_{\NS}}{2} \cdot C \cdot \tau_{\text{indirect}}
\end{equation}
\end{proof}

\textbf{Key Implication}: The more connected an entity, the harder it is to hide. Hub removal creates correlated tension across the entire adjacency layer.

% ============================================
% 3.5 TEMPORAL DYNAMICS
% ============================================
\subsection{Temporal Dynamics: The Velocity of Silence}

The preceding analysis is \emph{static}---it characterizes tension at a single time point. We now develop the \emph{dynamic} theory that tracks how tension evolves over time, enabling estimation of suppression age.

\subsubsection{The Tension Diffusion Equation}

\begin{theorem}[Tension Diffusion]
\label{thm:diffusion}
Informational tension evolves according to the graph diffusion equation:
\begin{equation}
    \boxed{
    \frac{\partial \tau}{\partial t} = -\alpha L \tau + \mathcal{S}(v, t)
    }
\end{equation}
where $L$ is the graph Laplacian, $\alpha > 0$ is the diffusivity constant, and $\mathcal{S}(v,t)$ is the source term encoding active suppression.
\end{theorem}

\begin{proof}[Proof Sketch]
At node $v_i$, tension flows to neighbors proportionally to local gradients:
\begin{equation}
    \frac{\partial \tau_i}{\partial t} = -\sum_{j \sim i} \alpha w_{ij}(\tau_i - \tau_j) + \mathcal{S}_i = -\alpha (L\tau)_i + \mathcal{S}_i
\end{equation}
This is the discrete analogue of the continuous heat equation $\partial_t u = \alpha \nabla^2 u$.
\end{proof}

The source term distinguishes two suppression regimes:
\begin{align}
    \mathcal{S}^{\text{passive}}(v, t) &= \tau_0(v) \cdot \delta(t - t_0) \quad \text{(one-time removal)} \\
    \mathcal{S}^{\text{active}}(v, t) &= \sigma(v) \cdot \mathbf{1}_{t \geq t_0} \quad \text{(continuous concealment)}
\end{align}

\subsubsection{Spectral Solution and Decay Timescales}

Let $0 = \lambda_0 < \lambda_1 \leq \cdots \leq \lambda_{n-1}$ be the eigenvalues of the graph Laplacian $L$ with orthonormal eigenvectors $\{\phi_k\}$.

\begin{theorem}[Spectral Decomposition]
\label{thm:spectral}
For passive suppression with initial condition $\tau(v, 0) = \tau_0(v)$:
\begin{equation}
    \boxed{
    \tau(v, t) = \sum_{k=0}^{n-1} \langle \tau_0, \phi_k \rangle \cdot e^{-\alpha \lambda_k t} \cdot \phi_k(v)
    }
\end{equation}
\end{theorem}

\begin{corollary}[Characteristic Timescales]
The tension dynamics exhibit two timescales:
\begin{align}
    T_{\text{local}} &= (\alpha \lambda_{n-1})^{-1} \quad \text{(fast local equilibration)} \\
    T_{\text{global}} &= (\alpha \lambda_1)^{-1} \quad \text{(slow global decay)}
\end{align}
For $t \gg T_{\text{local}}$, only the Fiedler mode survives: $\tau(v, t) \approx c_1 e^{-\alpha \lambda_1 t} \phi_1(v)$.
\end{corollary}

\subsubsection{The Velocity of Silence}

\begin{definition}[Velocity of Silence]
\label{def:velocity}
The \emph{Velocity of Silence} $\vsilence$ is the characteristic propagation speed of tension through the network:
\begin{equation}
    \boxed{
    \vsilence = \alpha \sqrt{\lambda_1} \cdot \bar{\ell}
    }
\end{equation}
where $\lambda_1$ is the Fiedler eigenvalue (algebraic connectivity) and $\bar{\ell}$ is the mean edge length.
\end{definition}

\textbf{Physical interpretation}: $\vsilence$ measures how fast the ``news'' of a suppression propagates. High connectivity ($\lambda_1$ large) implies fast propagation; sparse networks have slow propagation.

\begin{proposition}[Age Estimation]
\label{prop:age}
Given observed tension profile $\tau_{\text{obs}}(v)$ with radial extent $r$ from the suppression epicenter, the suppression age is approximately:
\begin{equation}
    t_{\text{supp}} \approx \frac{r}{\vsilence} = \frac{r}{\alpha \sqrt{\lambda_1} \cdot \bar{\ell}}
\end{equation}
\end{proposition}

This enables \emph{temporal forensics}: distinguishing recent suppressions (compact, high-magnitude tension) from old suppressions (diffuse, low-magnitude tension).

% ============================================
% 4. THE YHARIM LIMIT
% ============================================
\section{The Yharim Limit: Detectability Threshold}

Not all suppressions are detectable. We now establish the limits of detection.

\subsection{Signal-to-Noise Analysis}

\begin{definition}[Suppression Signal-to-Noise Ratio]
The SNR for node $\NS$ with degree $d_{\NS}$, in a graph with degree variance $\sigma_d^2$ and local clustering $C_{\NS}$, is:
\begin{equation}
    \text{SNR}(\NS) = \frac{d_{\NS}}{\sigma_d} \cdot \sqrt{C_{\NS}}
\end{equation}
\end{definition}

The SNR captures how much the suppression signal ``stands out'' from background variation. High-degree nodes in clustered regions produce strong signals; low-degree nodes in sparse regions produce weak signals indistinguishable from noise.

\subsection{The Detectability Threshold}

\begin{definition}[The Yharim Limit]
\label{def:yharim}
The \emph{Yharim Limit} $\Yharim$ is the fundamental threshold below which informational tension is indistinguishable from stochastic noise:
\begin{equation}
    \boxed{
    \Yharim \equiv \sqrt{2 \log(1/\alpha)}
    }
\end{equation}
where $\alpha$ is the desired false positive rate. Reference values:
\begin{center}
\renewcommand{\arraystretch}{1.2}
\begin{tabular}{ccc}
\toprule
\textbf{Significance} ($\alpha$) & \textbf{Yharim Limit} ($\Yharim$) & \textbf{Interpretation} \\
\midrule
0.05 & 2.45 & Standard threshold \\
0.01 & 3.03 & Stringent threshold \\
0.001 & 3.72 & High-confidence threshold \\
\bottomrule
\end{tabular}
\end{center}
\end{definition}

The Yharim Limit derives from scan statistic theory: when testing multiple regions for anomalies, the maximum of many Gaussian variables follows a Gumbel distribution with threshold $\sqrt{2\log n}$.

\begin{theorem}[Detectability Threshold]
\label{thm:detectability}
A suppressed node $\NS$ is statistically detectable if and only if:
\begin{equation}
    \boxed{\text{SNR}(\NS) > \Yharim}
\end{equation}
Equivalently, the minimum detectable degree is:
\begin{equation}
    d_{\min} = \frac{\sigma_d}{\sqrt{C_{\NS}}} \cdot \Yharim
\end{equation}
\end{theorem}

\begin{proof}[Proof Sketch]
The detection problem reduces to identifying a mean shift in a Gaussian variable. By the Neyman-Pearson lemma, the optimal test compares the likelihood ratio to a threshold determined by $\alpha$. The critical SNR follows from requiring power $> 1/2 + \epsilon$ at significance level $\alpha$.
\end{proof}

\textbf{Interpretation}: The Yharim Limit represents the detectability horizon---the threshold beyond which no statistical method can distinguish suppression from natural variation. Low-degree nodes in low-clustering regions are undetectable by any method.

\subsection{Validity of the Yharim Limit for Scale-Free Networks}

A critical question: does the Gaussian-derived threshold $\Yharim = \sqrt{2\log(1/\alpha)}$ apply to scale-free networks where degree distributions have infinite variance?

\begin{theorem}[JSD Boundedness Guarantees Applicability]
\label{thm:jsd-scalefree}
Let $\Graph$ be a scale-free network with degree distribution $P(k) \sim k^{-\gamma}$, $\gamma \in (2, 3)$. Although $\sigma_d^2 = \infty$, the JSD-based tension $\tau_{\text{JSD}}(v)$ satisfies conditions for extreme value theory, and the Yharim threshold remains valid.
\end{theorem}

\begin{proof}
The proof proceeds in three steps.

\textbf{Step 1: JSD is bounded.} By definition:
\begin{equation}
    \tau_{\text{JSD}}(v) = \Djsd(P_{\text{obs}} \| P_{\text{exp}}) \in [0, 1]
\end{equation}
Unlike KL divergence (which can be infinite), JSD is bounded regardless of the underlying degree distribution. This is the key property that ``tames'' heavy tails.

\textbf{Step 2: Bounded random variables have finite moments.} For any random variable $X \in [0, 1]$:
\begin{equation}
    \mathbb{E}[X^k] \leq 1 \quad \text{for all } k \geq 1
\end{equation}
Therefore $\tau_{\text{JSD}}$ has finite variance $\sigma_\tau^2 < 1$, even when computed on scale-free networks.

\textbf{Step 3: Extreme value convergence.} By the Fisher-Tippett-Gnedenko theorem, if $\{X_i\}$ are i.i.d.\ with finite variance, then $\max_i X_i$ converges to a Gumbel distribution. For $n$ nodes tested simultaneously:
\begin{equation}
    \Prob\left( \max_i \tau_i > \mu + \sigma \sqrt{2 \log n} \right) \to e^{-1} \approx 0.368
\end{equation}
Setting the false positive rate to $\alpha$ yields threshold $\sqrt{2\log(1/\alpha)}$.

\textbf{Conclusion}: The Yharim Limit applies to JSD-based tension in scale-free networks because JSD boundedness ensures finite variance, which guarantees Gumbel-domain convergence.
\end{proof}

\begin{remark}[Why KL Divergence Fails]
If we used raw KL divergence instead of JSD, the threshold would \emph{not} be valid for scale-free networks. KL divergence is unbounded and can produce heavy-tailed tension distributions that violate Gumbel-domain assumptions. The choice of JSD is not merely ``robust''---it is \emph{mathematically necessary} for the Yharim Limit to hold in real networks.
\end{remark}

% ============================================
% 5. THE SNIPER EFFECT
% ============================================
\section{The Sniper Effect: Detecting Low-Degree Suppressions}

A critical and perhaps counterintuitive prediction of ITT: suppressed nodes can leave detectable entropic traces \emph{regardless of their own degree}. The key is not the connectivity of the suppressed node, but the connectivity of its neighbors.

\subsection{Theoretical Foundation}

\begin{definition}[The Sniper Effect]
\label{def:sniper}
Node $v$ with low degree $d(v) \ll \bar{d}$ exhibits the \emph{Sniper Effect} if:
\begin{equation}
    \tau(v) > \tau_{\text{threshold}} \quad \text{despite} \quad d(v) < d_{\min}
\end{equation}
This occurs when $v$'s neighbors are well-connected nodes that lost a common connection---the suppressed ``target'' that the low-degree ``sniper'' was pointing at.
\end{definition}

The name evokes the analogy: a sniper has few connections (low degree) but impacts high-value targets. When the sniper is removed, the targets' behavior reveals the absence.

\begin{theorem}[Hermit-Suppression Discriminant]
\label{thm:hermit}
Let $\Sigma$ be a low-degree void region. Define the \emph{neighbor connectivity score}:
\begin{equation}
    \text{NCS}(\Sigma) = \frac{1}{|\partial\Sigma|} \sum_{v \in \partial\Sigma} d(v)
\end{equation}
where $\partial\Sigma$ is the boundary of $\Sigma$. Then:
\begin{itemize}
    \item $\text{NCS}(\Sigma) \approx \bar{d}$ $\Rightarrow$ Natural hermit (low tension expected)
    \item $\text{NCS}(\Sigma) \gg \bar{d}$ $\Rightarrow$ Sniper suppression (high tension expected)
\end{itemize}
\end{theorem}

The discriminant asks: \emph{``Are the nodes surrounding this void unusually well-connected?''} If yes, suppression is indicated regardless of the void's apparent size.

\subsection{Empirical Validation: The Marak Incident}

\begin{example}[The Marak Incident]
\label{ex:marak}
In January 2022, developer ``Marak Squires'' was banned from GitHub after deliberately corrupting the \texttt{faker.js} and \texttt{colors.js} packages, affecting thousands of dependent projects.

\textbf{Traditional analysis}: Marak had sampled degree $d = 2$ in the co-contribution network, placing him in the theoretically \emph{undetectable} zone ($P < 0.1$) according to degree-based metrics.

\textbf{ITT analysis}: Despite low degree, Marak's local tension was $\tau = 0.31$, exceeding both the detection threshold ($\tau_{\text{crit}} = 0.2$) and the tension of high-degree ecological hubs ($\tau = 0.29$ for Claytonia with $d = 66$).

\textbf{Conclusion}: The Sniper Effect validated---despite topological invisibility ($d < d_{\min}$), the suppression was detectable via elevated JSD in the adjacency layer. Degree-based methods yield false negative; divergence-based methods yield true positive.
\end{example}

% ============================================
% 6. EMPIRICAL VALIDATION
% ============================================
\section{Empirical Validation: Cross-Domain Results}

To validate the theoretical predictions, we applied ITT to two radically different domains:

\begin{enumerate}
    \item \textbf{Software Ecosystems}: GitHub co-contribution network around the \texttt{faker.js} package (December 29, 2021 -- January 6, 2022), analyzing $N=62$ developers
    
    \item \textbf{Ecological Networks}: Plant-pollinator mutualistic network from Burkle et al.\ (2013), with 78 nodes and $\sim$500 interactions, simulating extinction of \textit{Claytonia virginica}
\end{enumerate}

\subsection{Quantitative Results}

\begin{table}[htbp]
\centering
\caption{Cross-Domain Validation Results}
\label{tab:validation}
\renewcommand{\arraystretch}{1.3}
\begin{tabular}{llcccc}
\toprule
\textbf{Domain} & \textbf{Suppressed Node} & \textbf{$d$} & \textbf{$\tau$} & \textbf{$\Omega$} & \textbf{Verdict} \\
\midrule
Ecology & \textit{Claytonia virginica} & 66 & 0.293 & 19.32 & \textsc{High Tension} \\
Software & \texttt{user:Marak} & 2 & 0.311 & 0.62 & \textsc{High Tension} \\
\bottomrule
\end{tabular}

\vspace{0.5em}
\footnotesize
\textit{Note}: Both exhibit $\tau > 0.2$ (threshold) despite 33$\times$ difference in degree. Tension computed using JSD; thresholding uses MAD for robustness to power-law outliers. Data sources: GitHub Archive and Burkle et al.\ (2013).
\end{table}

\subsection{Validation of Theoretical Predictions}

\begin{enumerate}
    \item \textbf{Axiom 2 (Deformation) Validated}: Both domains show $\tau > 0.2$, confirming that suppression creates measurable structural deformation
    
    \item \textbf{Theorem \ref{thm:scaling} (Superlinear Scaling) Validated}: Hub concealment cost ($\Omega \approx 19.3$) exceeds peripheral node cost ($\Omega \approx 0.2$) by approximately 100$\times$, consistent with superlinear scaling
    
    \item \textbf{The Sniper Effect Confirmed}: Marak detected via high $\tau$ despite $d = 2 < d_{\min}$, validating that low-degree suppressions leave detectable traces when neighbors are well-connected
    
    \item \textbf{Domain Equivalence}: The normalized concealment cost $\Omega/d$ is strikingly similar:
    \begin{itemize}
        \item Ecology: $19.32 / 66 = 0.29$
        \item Software: $0.62 / 2 = 0.31$
    \end{itemize}
    This validates the domain equivalence theorem: the same mathematical structure governs structurally isomorphic suppression events across domains.
\end{enumerate}

\subsection{Methodology Notes}

\begin{itemize}
    \item \textbf{Robust Estimation}: All validations use JSD (bounded $[0, 1]$) rather than KL divergence, and MAD-based thresholds ($\tau_{\text{thresh}} \approx 1.0$) rather than standard deviation
    
    \item \textbf{Reproducibility}: Data from GitHub Archive (public) and Dryad repository (Burkle et al., 2013). Analysis code available at \texttt{github.com/MatheusGrego/itt-poc}
    
    \item \textbf{Limitations}: The software validation uses a single incident; future work should validate across multiple events (e.g., \texttt{event-stream}, \texttt{ua-parser-js})
\end{itemize}

% ============================================
% 7. DOMAIN EQUIVALENCE
% ============================================
\section{Domain Equivalence Theorem}

\begin{theorem}[Domain Equivalence]
\label{thm:domain-equiv}
Let $\mathfrak{S}_1$, $\mathfrak{S}_2$ be information flow systems from different domains. If there exists a bijection $\psi: \mathcal{E}_1 \to \mathcal{E}_2$ preserving structural isomorphism, flow proportionality, and capacity proportionality, then:
\begin{equation}
    \tau_{\mathfrak{S}_1}(v; \mathcal{S}_e) = \tau_{\mathfrak{S}_2}(\psi(v); \mathcal{S}_{\psi(e)})
\end{equation}
\begin{equation}
    \Omega_{\mathfrak{S}_1}(\mathcal{S}_e) = \Omega_{\mathfrak{S}_2}(\mathcal{S}_{\psi(e)})
\end{equation}
Tension and concealment cost are \emph{invariant} under domain transformation.
\end{theorem}

\textbf{Implication}: ITT is a \emph{domain-agnostic} framework. Applications include:
\begin{itemize}[noitemsep]
    \item \textbf{Cybersecurity}: Detecting removed malicious packages in supply chains
    \item \textbf{Finance}: Identifying dissolved shell companies in transaction networks
    \item \textbf{Social}: Detecting censored accounts in social networks
    \item \textbf{Biology}: Inferring extinct species from ecosystem structure
\end{itemize}

% ============================================
% 8. CONCLUSION
% ============================================
\section{Conclusion}

Informational Tension Theory provides a mathematical framework for detecting and dating latent entities through structural analysis. The key contributions are:

\begin{enumerate}
    \item \textbf{Three Axioms}: Conservation, Deformation, and Structural Cost ground the theory
    \item \textbf{Informational Tension} ($\tau$): A local, JSD-based anomaly measure with proven validity in scale-free networks
    \item \textbf{Concealment Cost} ($\Omega$): A global measure with superlinear scaling $\Omega \sim d^{1+\alpha}$
    \item \textbf{Temporal Dynamics}: The diffusion equation $\partial_t \tau = -\alpha L\tau + \mathcal{S}$ and the \emph{Velocity of Silence} $\vsilence = \alpha\sqrt{\lambda_1}\bar{\ell}$ enable suppression age estimation
    \item \textbf{The Yharim Limit} ($\Yharim$): The detection threshold $\sqrt{2\log(1/\alpha)}$, proven valid for JSD in scale-free networks
    \item \textbf{The Sniper Effect}: Low-degree nodes detectable via adjacency layer divergence
    \item \textbf{Empirical Validation}: Cross-domain confirmation (Software $\leftrightarrow$ Ecology)
\end{enumerate}

The framework provides both \emph{static} detection (where is suppression?) and \emph{dynamic} forensics (when did suppression occur?). Open problems include empirical validation of temporal predictions and adversarial robustness analysis.

% ============================================
% REFERENCES
% ============================================
\section*{References}

\begin{enumerate}[label={[\arabic*]}, leftmargin=*, noitemsep]
    \item Burkle, L.A., Marlin, J.C., \& Knight, T.M. (2013). Plant-Pollinator Interactions over 120 Years. \textit{Science}, 339(6127), 1611--1615.
    
    \item Barabási, A.L. \& Albert, R. (1999). Emergence of Scaling in Random Networks. \textit{Science}, 286(5439), 509--512.
    
    \item Lin, J. (1991). Divergence measures based on the Shannon entropy. \textit{IEEE Trans. Inf. Theory}, 37(1), 145--151.
    
    \item Chung, F.R.K. (1997). \textit{Spectral Graph Theory}. CBMS Regional Conference Series, AMS.
    
    \item Ollivier, Y. (2009). Ricci curvature of Markov chains on metric spaces. \textit{J. Functional Analysis}, 256(3), 810--864.
    
    \item Roy, A. et al. (2023). GAD-NR: Graph Anomaly Detection via Neighborhood Reconstruction. \textit{ACM Web Conference}.
    
    \item Hampel, F.R. (1974). The Influence Curve and Its Role in Robust Estimation. \textit{JASA}, 69(346), 383--393.
    
    \item Fisher, R.A. \& Tippett, L.H.C. (1928). Limiting forms of the frequency distribution of the largest or smallest member of a sample. \textit{Proc. Cambridge Phil. Soc.}, 24, 180--190.
\end{enumerate}

\vspace{1em}
\hrule
\vspace{0.5em}
\footnotesize
\textbf{Full Paper}: Extended version (64 pages) with complete proofs, geometric analysis, temporal dynamics, and additional validation available from the author.

\textbf{Code}: Proof-of-concept implementation available at \texttt{github.com/MatheusGrego/itt-poc}

\textbf{Data}: GitHub Archive (Dec 2021--Jan 2022); Dryad repository (Burkle et al., 2013)

\end{document}
